% Part: computability
% Chapter: recursive-functions
% Section: introduction

\documentclass[../../../include/open-logic-section]{subfiles}

\begin{document}

\olfileid{cmp}{rec}{int}
\olsection{పరిచయం}

గణనీయతకు ఒక గణిత సిద్ధాంతాన్ని అభివృద్ధి చేయాలంటే, అన్నిటికంటే
ముందుగా గణనీయతకు ఒక \emph{నమూనా}ను రూపొందించాలి. ఇప్పుడు గణనీయతను
కంప్యూటర్లు చేసే పని రకంగా మనం భావిస్తాం; కంప్యూటర్లు సంకేతాలతో
పనిచేస్తాయి. అయితే గణనీయతా సిద్ధాంతాలు మొదట అభివృద్ధి చెందినప్పుడు,
గణనకు ఆదర్శ ఉదాహరణ \emph{సంఖ్యాత్మక} గణన. గణితవేత్తలు ఎప్పటినుంచో
గణించగల సంఖ్యా-సిద్ధాంత ప్రమేయాలపై, అంటే $f\colon \Nat^n \to
\Nat$ రూపపు ప్రమేయాలపై, ఆసక్తి చూపారు. అందువల్ల గణనీయతా సిద్ధాంతం
ఆరంభంలో అలాంటి ప్రమేయాలనే అధ్యయనం చేయడంలో ఆశ్చర్యం లేదు. కూడిక,
గుణకారం, సహజ సంఖ్యల ఘాతాంకనం వంటి సుపరిచిత గణనీయ సంఖ్యాత్మక
ప్రమేయాలన్నిటికీ ఒక ఆసక్తికర లక్షణం ఉంది: వాటిని \emph{పునరావృత్తంగా}
నిర్వచించవచ్చు. కనుక పునరావృత్త నిర్వచనాల ఆధారంగా \emph{గణనీయ
ప్రమేయం}కు ఒక సాధారణ నిర్వచనం ఇవ్వాలని ప్రయత్నించడం సహజమే.
సంఖ్యా-సిద్ధాంత ప్రమేయాలను పునరావృత్తంగా నిర్వచించగల అనేక పద్ధతుల్లో,
ముఖ్యమైనదిగా మారే ఒక ప్రత్యేకంగా సరళమైన నమూనా ఉంది: దానినే
\emph{ఆదిమ పునరావృత్తి} అంటారు.

గణనీయ ప్రమేయాలతో పాటు గణనీయ సమితులు, సంబంధాలపైనా మనకు ఆసక్తి
ఉండవచ్చు. ఇచ్చిన సంఖ్య ఆ సమితిలో \tetoken{ఒక మూలకం}{!!a{element}} అవునో కాదో సమాధానాన్ని
గణించగలిగితే ఆ సమితి గణనీయమైనది; ఇచ్చిన ట్యుపుల్
$\tuple{n_1, \dots, n_k}$ ఆ సంబంధంలో \tetoken{ఒక మూలకం}{!!a{element}} అవునో కాదో
గణించగలిగితే, అప్పుడే ఆ సంబంధం గణనీయమైనది. ఒక సమితి లేదా సంబంధం
యొక్క \emph{లక్షణ ప్రమేయం}ను పరిగణించడం ద్వారా, గణనీయ సమితులు మరియు
సంబంధాల చర్చను గణనీయ ప్రమేయాల చర్చలోనే చేర్చవచ్చు. అందువల్ల ఆదిమ
పునరావృత్త సంబంధాలను కూడా నిర్వచించవచ్చు; ఉదాహరణకు, ``$n$, $m$ను
శేషంలేకుండా భాగిస్తుంది'' అనే సంబంధం ఆదిమ పునరావృత్త సంబంధం.

అయితే కేవలం ఆదిమ పునరావృత్తిని ఉపయోగించి నిర్వచించగల ఆదిమ
పునరావృత్త ప్రమేయాలు మాత్రమే గణనీయ సంఖ్యా-సిద్ధాంత ప్రమేయాలు కావు.
ఆదిమ పునరావృత్తికి అనేక సాధారణీకరణలను పరిశీలించారు; వాటిలో అత్యంత
శక్తివంతమైన, విస్తృతంగా అంగీకరించిన అదనపు గణన పద్ధతి అపరిమిత
అన్వేషణ. ఇది \emph{పాక్షిక పునరావృత్త ప్రమేయాల} నిర్వచనానికీ,
దానికి సంబంధించిన \emph{సామాన్య పునరావృత్త ప్రమేయాల} నిర్వచనానికీ
దారితీస్తుంది. సామాన్య పునరావృత్త ప్రమేయాలు గణనీయమైనవి,
సర్వనిర్వచితమైనవి; సర్వనిర్వచితంగా ఉన్న పాక్షిక పునరావృత్త
ప్రమేయాలనే ఈ నిర్వచనం సరిగ్గా లక్షణీకరిస్తుంది. పునరావృత్త ప్రమేయాలు
ప్రతి ఇతర గణనా నమూనానూ (ట్యూరింగ్ యంత్రాలు, లాంబ్డా కలనశాస్త్రం
మొదలైనవి) అనుకరించగలవు; అందువల్ల అవి అంగీకృత గణనా నమూనాల్లో ఒకటి.


\end{document}
